% ---- pg-569 （书页 557） ----

\begin{mvfigure}
\mvinc{pg-569}{0.192,0.121,0.818,0.537}
\phantomsection\label{fig:12.20}\mvcap{图 12.20\quad 由非度量多维标度变换得出的大学的二维表示}
\end{mvfigure}

\phantomsection\label{eq:12-27}
\begin{equation}
E=\left[\sum_{i<k}\sum\left(d_{ik}^{(p)}-d_{ik}^{(q)}\right)^2\Big/d_{ik}^{(p)}\right]
\left[\sum_{i<k}\sum d_{ik}^{(p)}\right]^{-1}
\tag{12-27}
\end{equation}

达到最小来找到 $q<p$ 维表示. 在这种替代方法中, 将 $p$ 维和 $q$ 维欧氏距离进行直接比较. 通过使 $E$ 最小化来得到低维表示的方法称为非线性映射.

任一个低维表示的最后拟合优度均可用最小张树来表示(关于这个问题的进一步讨论见文献[16]).

\section{对应分析}
\phantomsection\label{sec:12.7}

对应分析是一种将频数或计数表中的各种联系用图来表示的方法, 由弗伦奇(French)所提出. 我们将集中讨论二维的频数表, 常称为列联表. 若列联表有 $I$ 行和 $J$ 列, 则由对应分析得出的图包含两组点, 一组有 $I$ 个点, 对应 $I$ 行; 另一组有 $J$ 个点, 对应 $J$ 列. 各点的位置反映各种联系.

相互接近的行点表明相应的行有相似的轮廓(条件分布). 相互接近的列点表明相应的列有相似的轮廓(条件分布). 最后, 行点和列点相互接近表明相应的组合发生的频数会高于独立模型(行类别与列类别之间没有关系的模型)下的期望值.

对应分析的输出包含数据的“最佳”二维表示、被标绘各点的坐标值以及每个维中所保留的信息量的一个度量(称为惯性).
